\documentclass[a4paper]{article}
\usepackage{amsfonts}
\usepackage{amsmath}
\usepackage{amssymb}
\usepackage{graphicx}     

\newcommand{\dint}{\displaystyle\int}
\newcommand{\llim}{\lim\limits}

\begin{document}
  
\noindent \large Auteur: Oubayy Ahale          \\
\noindent \large Studierichting: Informatica   \\
\noindent \large Oefeningenreeks 2E nr. 26.16  \\

\medskip

\normalsize

$ 2 \cdot (\log_2 (x-1) - \log_4 (x^2+1)) = 1 - \log_4 25 $\\

$ \Leftrightarrow 2 \cdot \left( \log_2 (x-1) - \dfrac{\log_2 (x^2+1)}{\log_2 4} \right) = 1 - \dfrac{\log_2 25}{\log_2 4} $\\

$ \Leftrightarrow 2 \cdot \left( \log_2 (x-1) - \dfrac{\log_2 (x^2+1)}{2} \right) = 1 - \dfrac{\log_2 25}{2} $\\

$ \Leftrightarrow 2 \cdot \log_2 (x-1) - \log_2 (x^2+1) = 1 - \dfrac{\log_2 25}{2} $\\

$ \Leftrightarrow \log_2 (x-1)^2 - \log_2 (x^2+1) = 1 - \dfrac{\log_2 25}{2} $\\

$ \Leftrightarrow \log_2 \left( \dfrac{(x-1)^2}{x^2+1} \right) = 1 - \dfrac{\log_2 25}{2} $\\

$ \Leftrightarrow \log_2 \left( \dfrac{(x-1)^2}{x^2+1} \right) = \log_2 2 - \dfrac{2 \log_2 5}{2} $\\

$ \Leftrightarrow \log_2 \left( \dfrac{(x-1)^2}{x^2+1} \right) = \log_2 2 - \log_2 5 $\\

$ \Leftrightarrow \log_2 \left( \dfrac{(x-1)^2}{x^2+1} \right) = \log_2 \dfrac{2}{5} $\\

$ \Leftrightarrow \dfrac{(x-1)^2}{x^2+1} = \dfrac{2}{5} $\\

$ \Leftrightarrow (x-1)^2 = \dfrac{2}{5} (x^2+1) $\\

$ \Leftrightarrow x^2 - 2x + 1 = \dfrac{2}{5} x^2 + \dfrac{2}{5} $\\

$ \Leftrightarrow 5x^2 - 10x + 5 = 2x^2 + 2 $\\

$ \Leftrightarrow 3x^2 - 10x + 3 = 0 $\\

$ D = b^2 - 4ac $\\

$ = 100 - 4 \cdot 3 \cdot 3 $\\

$ = 100 - 36 $\\

$ = 64 $\\

$ x_1 = \dfrac{-b + \sqrt{D}}{2a} = \dfrac{10 + 8}{6} = \dfrac{18}{6} = 3 $\\

$ x_2 = \dfrac{-b - \sqrt{D}}{2a} = \dfrac{10 - 8}{6} = \dfrac{2}{6} = \dfrac{1}{3} $\\

$ V = \{ 3 \} $\\

\begin{thebibliography}{999}
%\bibitem{a} Referentie
\end{thebibliography}
\end{document} 